% !TeX program = xelatex
% 使用 XeLaTeX；Fandol 字体随 TeX Live 提供，无需系统字体。
\documentclass[UTF8,fontset=fandol,a4paper,zihao=-4]{ctexart}
\usepackage{uestcdesign}
% 修改以下字段即可开始写作；含逗号的值请放在花括号内。
\uestcsetup{
theme=proposal,
 type=研究计划,english={RESEARCH PROPOSAL},
 title={低功耗感知系统设计},subtitle={问题定义、技术路线与阶段安排},
 author=张同学,student-id=202600000000,
 school=信息与通信工程学院,course={科研训练项目},
 teacher=教师姓名,date={2026 年 10 月 13 日},short-title={低功耗感知}
}
\begin{document}
\uestccover
\begin{summarybox}[项目摘要]
\textbf{研究问题：}如何在有限电池容量下持续采集环境信息？\\
\textbf{拟议方法：}分析系统能耗，联合设计采样周期与通信策略。\\
\textbf{预期产出：}可复现实验原型、测量数据、研究报告。\\
\textbf{性质：}这是研究计划书的示例框架，目标均需通过实际实验验证。
\end{summarybox}
\section{研究背景与问题}
感知节点在采集、处理、传输和休眠阶段具有不同功耗。单纯降低某一模块的功率，未必降低完整任务的能量消耗。因此需要明确工作负载和测量边界。
\section{研究目标与评价指标}
\noindent\begin{tabularx}{\linewidth}{@{}lXX@{}}\toprule
目标 & 评价指标 & 验证方式\\\midrule
降低能耗 & 单次任务能量 & 电流与电压同步测量\\
保持可靠性 & 数据完整率 & 丢包与异常记录\\
控制响应时间 & 端到端延迟 & 统一时钟下记录时间戳\\\bottomrule
\end{tabularx}
\section{技术路线}
\begin{enumerate}
\item \textbf{建立基线：}固定采样周期，记录各工作阶段的时间与功率。
\item \textbf{定位瓶颈：}根据能量积分比较采集、计算、通信的贡献。
\item \textbf{提出策略：}设计自适应采样与批量传输方案。
\item \textbf{验证假设：}在相同负载和环境条件下比较基线与改进方案。
\end{enumerate}
\begin{equation}
E_{\mathrm{task}}=\int_0^{T} U(t)I(t)\,\mathrm{d}t.
\end{equation}
应明确测量点是在电池侧还是模块侧，以及是否包含外围设备能耗。
\clearpage
\section{阶段计划与里程碑}
\begin{table}[htbp]\centering
\caption{八周研究安排示例}
\begin{tabularx}{\linewidth}{@{}lXX@{}}\toprule
阶段 & 工作内容 & 可检查的产出\\\midrule
第 1--2 周 & 梳理文献，明确测量口径 & 相关工作表、问题定义\\
第 3--4 周 & 搭建系统并测量基线 & 硬件清单、基线数据\\
第 5--6 周 & 实现策略，进行对比实验 & 代码、实验日志\\
第 7--8 周 & 分析结果，整理复现材料 & 报告、数据与演示\\\bottomrule
\end{tabularx}
\end{table}
\section{风险与备选方案}
\subsection{设备交付或测量条件受限}
先以开发板和合成负载验证流程，记录其与最终实验条件的差别；替代实验的结论不直接外推到目标系统。
\subsection{改进效果不明显}
保留无效实验并分析瓶颈，检查策略引入的额外开销，调整研究假设或缩小适用范围。
\section{资源与预期成果}
需要具备电流测量条件的供电设备、开发板和数据存储空间。成果包括系统说明、源代码、数据字典及复现步骤。
\begin{summarybox}[阶段验收]
每个里程碑同时检查“是否完成”和“是否有证据”。当条件发生变化时，记录调整原因、影响范围和新的交付时间。
\end{summarybox}
\end{document}
